\documentclass[11pt,a4paper]{article}
\usepackage[utf8]{inputenc}
\usepackage[T1]{fontenc}
\usepackage[romanian]{babel}
\usepackage[margin=2.5cm]{geometry}
\usepackage{graphicx} % Required for inserting images
\usepackage{xcolor}
\usepackage{listings}
\usepackage{booktabs}
\usepackage{tabularx}
\usepackage{enumitem}
\usepackage{hyperref}

\hypersetup{colorlinks=true, linkcolor=blue!50!black, urlcolor=blue!60!black}

\definecolor{codebg}{RGB}{246,246,246}
\lstset{
  basicstyle=\ttfamily\small,
  backgroundcolor=\color{codebg},
  frame=single,
  rulecolor=\color{gray!40},
  breaklines=true,
  columns=fullflexible,
  keepspaces=true,
  showstringspaces=false,
  numbers=left,
  numberstyle=\tiny\color{gray},
  xleftmargin=1.5em
}

\newcommand{\done}{\textcolor{green!50!black}{\textbf{realizat}}}
\newcommand{\wip}{\textcolor{blue!60!black}{\textbf{în lucru}}}
\newcommand{\todo}{\textcolor{orange!80!black}{\textbf{planificat}}}

\title{SecondChancePaws}
\author{Alexandra-Mihaela Bontidean}
\date{November 2026}

\begin{document}

\maketitle
\tableofcontents
\newpage

\section{Introducere}

\textbf{Second Chance Paws} este o aplicație web pentru sprijinirea adopției animalelor
fără stăpân și a centrelor care le adăpostesc. Aplicația este un proiect personal,
dezvoltat în scop de învățare și portofoliu, cu orientare spre folosul comunității.

Platforma pune în legătură trei categorii de persoane:
\begin{itemize}
  \item persoane care doresc să \textbf{adopte} un animal;
  \item persoane care \textbf{nu mai pot avea grijă} de animalul lor sau care
        \textbf{au găsit} un animal pe stradă și vor să-l ducă la un centru;
  \item \textbf{centrele de adopție} (adăposturi, ONG-uri), care își gestionează
        animalele și cererile primite.
\end{itemize}

\subsection{Obiective}
\begin{itemize}
  \item afișarea animalelor disponibile pentru adopție, cu filtrare și paginare;
  \item o hartă interactivă cu centrele, datele de contact și distanța până la ele;
  \item fluxuri separate pentru adopție, predarea unui animal și raportarea unui animal găsit;
  \item integrarea unui asistent bazat pe inteligență artificială;
  \item rularea întregii aplicații în containere Docker.
\end{itemize}

\subsection{Date folosite}
Aplicația folosește exclusiv \textbf{date demonstrative fictive}: centre cu nume
inventate, amplasate pe harta orașului Cluj-Napoca, numere de telefon evident false
(de forma \texttt{0700 000 001}) și animale generate automat. Interfața afișează un
mesaj care precizează că datele sunt fictive.

\section{Tehnologii}

\begin{center}
\begin{tabularx}{\textwidth}{@{}lX@{}}
\toprule
\textbf{Componentă} & \textbf{Tehnologie} \\
\midrule
Limbaj backend      & Java 21 \\
Framework backend   & Spring Boot 4.1 (Spring Web, Spring Data JPA, Spring Security) \\
ORM                 & Hibernate \\
Bază de date        & PostgreSQL 17 \\
Migrări bază de date & Flyway \\
Autentificare       & JWT (JSON Web Token) \\
Documentare API     & OpenAPI / Swagger UI (springdoc) \\
Frontend            & React + TypeScript (Vite) \\
Hartă               & Leaflet + OpenStreetMap \\
Inteligență artificială & Claude API (Anthropic) \\
Containerizare      & Docker, Docker Compose \\
Administrare BD     & pgAdmin 4 \\
\bottomrule
\end{tabularx}
\end{center}

\subsection{PostgreSQL, Flyway și Hibernate}
Cele trei tehnologii nu sunt alternative, ci se completează:
\begin{itemize}
  \item \textbf{PostgreSQL} este sistemul de gestiune a bazei de date, unde sunt stocate datele.
  \item \textbf{Flyway} gestionează structura bazei de date prin scripturi SQL versionate
        (\texttt{V1\_\_init.sql}, \texttt{V2\_\_demo\_data.sql} etc.), rulate automat
        la pornirea aplicației, în ordine.
  \item \textbf{Hibernate} mapează tabelele pe clase Java (entități). Este configurat cu
        \texttt{ddl-auto=validate}: verifică faptul că entitățile corespund schemei,
        fără a o modifica. Schema este controlată exclusiv de Flyway.
\end{itemize}

\section{Arhitectură}

Aplicația are o arhitectură client--server, cu trei componente principale, fiecare
rulând în propriul container:

\begin{enumerate}
  \item \textbf{Frontend} (React) -- interfața din browser; comunică cu backend-ul prin HTTP/JSON.
  \item \textbf{Backend} (Spring Boot) -- API REST, logica aplicației, securitate,
        comunicarea cu serviciul de inteligență artificială.
  \item \textbf{Baza de date} (PostgreSQL) -- stocarea persistentă a datelor.
\end{enumerate}

Backend-ul este organizat pe straturi:
\texttt{controller} (endpoint-uri REST) $\rightarrow$ \texttt{service} (logica aplicației)
$\rightarrow$ \texttt{repository} (acces la date prin Spring Data JPA) $\rightarrow$ baza de date.

\section{Roluri și utilizatori}

\begin{center}
\begin{tabularx}{\textwidth}{@{}lX@{}}
\toprule
\textbf{Rol} & \textbf{Permisiuni} \\
\midrule
\texttt{USER}          & caută animale, trimite cereri de adopție, predă un animal,
                         raportează un animal găsit \\
\texttt{SHELTER\_STAFF} & gestionează animalele centrului, aprobă sau respinge cereri \\
\texttt{ADMIN}         & aprobă centrele nou înscrise, moderează conținutul \\
\bottomrule
\end{tabularx}
\end{center}

Un singur tip de cont este folosit pentru public, deoarece aceeași persoană poate
atât să adopte, cât și să predea sau să raporteze un animal. Centrele își pot crea
singure cont, iar acesta devine activ după aprobarea unui administrator.

Pentru testare există câte un utilizator demonstrativ pentru fiecare rol.

\section{Modelul de date (\todo)}

Entitățile principale planificate:
\begin{itemize}
  \item \textbf{User} -- datele contului și rolul;
  \item \textbf{Shelter} -- nume, adresă, latitudine, longitudine, telefon, program,
        capacitate, tipurile de animale acceptate, status de aprobare;
  \item \textbf{Animal} -- specie, rasă, vârstă, sex, mărime, sterilizat/vaccinat,
        fotografii, status (\texttt{AVAILABLE}, \texttt{RESERVED}, \texttt{ADOPTED}),
        centrul de care aparține;
  \item \textbf{AdoptionRequest} -- utilizator, animal, mesaj, status;
  \item \textbf{IntakeRequest} -- tip (\texttt{OWNER\_SURRENDER} sau \texttt{FOUND\_STRAY}),
        fotografii, locația, starea animalului, centrul ales, status;
  \item \textbf{Favorite} -- animalele salvate de un utilizator.
\end{itemize}

\section{Funcționalități}

\subsection{Listare animale cu paginare și filtre (\todo)}
Paginarea folosește interfețele \texttt{Pageable} și \texttt{Page<T>} din Spring Data.
Filtrele dinamice (specie, vârstă, mărime, compatibilitate cu copii, sterilizare) sunt
implementate cu \texttt{Specification}. Rezultatele pot fi sortate și după distanța
față de utilizator.

\subsection{Harta centrelor (\todo)}
\begin{itemize}
  \item harta este afișată cu Leaflet și OpenStreetMap (gratuit, fără cheie API);
  \item locația utilizatorului este obținută prin API-ul de geolocalizare al browserului;
  \item distanța în linie dreaptă se calculează cu formula Haversine:
\end{itemize}
\[
d = 2R \arcsin\left(\sqrt{\sin^2\frac{\varphi_2-\varphi_1}{2}
   + \cos\varphi_1 \cos\varphi_2 \sin^2\frac{\lambda_2-\lambda_1}{2}}\right)
\]
unde $R \approx 6371$ km este raza Pământului, $\varphi$ latitudinea și $\lambda$ longitudinea.
Fiecare centru afișează numele, telefonul, programul și distanța, plus un buton de navigare.

\subsection{Fluxul de adopție (\todo)}
Utilizatorul trimite o cerere pentru un animal; personalul centrului o aprobă sau o
respinge, iar statusul animalului se actualizează corespunzător.

\subsection{Predarea sau raportarea unui animal (\todo)}
Utilizatorul completează un formular (tip, fotografii, locație, descriere) și alege un
centru de pe hartă. Centrul primește cererea și o gestionează.

\subsection{Funcționalități pentru comunitate (\todo)}
\begin{itemize}
  \item anunțuri pentru animale pierdute și găsite;
  \item găzduire temporară (\emph{foster});
  \item lista de nevoi ale centrelor (hrană, pături, voluntari);
  \item evenimente și povești de succes ale adopțiilor.
\end{itemize}

\section{Integrarea inteligenței artificiale (\todo)}
Apelurile către Claude API se fac exclusiv din backend, cheia API nefiind expusă în
frontend. Scenarii planificate:
\begin{enumerate}
  \item \textbf{Potrivire la adopție} -- pe baza unui chestionar (locuință, copii, timp
        disponibil, alte animale), asistentul recomandă animale și explică alegerea;
  \item \textbf{Triajul unui animal găsit} -- din fotografie și descriere se estimează
        specia, vârsta și urgența, se oferă primii pași și se sugerează un centru;
  \item \textbf{Generarea descrierilor} -- profiluri de animale generate pentru personalul centrelor;
  \item \textbf{Asistent de chat} -- întrebări despre îngrijire și procesul de adopție.
\end{enumerate}

\section{Jurnal de implementare}

\subsection{Pasul 1: Baza de date în Docker (\done)}

PostgreSQL nu este instalat local, ci rulează într-un container Docker. Astfel mediul de
dezvoltare este reproductibil și poate fi pornit sau oprit cu o singură comandă.

Structura proiectului după acest pas:
\begin{lstlisting}
second-chance-paws/
|-- docker-compose.yml
|-- .env
|-- .gitignore
|-- .env.example
|-- README.md
|-- pgadmin/
|   `-- servers.json
`-- docs/
    `-- documentatie.tex
\end{lstlisting}

\paragraph{Fișierul \texttt{docker-compose.yml}} definește două servicii:
\begin{lstlisting}
services:
  db:
    image: postgres:17-alpine
    container_name: scp-db
    environment:
      POSTGRES_DB: ${POSTGRES_DB}
      POSTGRES_USER: ${POSTGRES_USER}
      POSTGRES_PASSWORD: ${POSTGRES_PASSWORD}
    ports:
      - "${POSTGRES_PORT:-5433}:5432"
    volumes:
      - pgdata:/var/lib/postgresql/data
    healthcheck:
      test: ["CMD-SHELL", "pg_isready -U ${POSTGRES_USER} -d ${POSTGRES_DB}"]
      interval: 5s
      timeout: 3s
      retries: 10

  pgadmin:
    image: dpage/pgadmin4:9
    container_name: scp-pgadmin
    environment:
      PGADMIN_DEFAULT_EMAIL: admin@demo.com
      PGADMIN_DEFAULT_PASSWORD: admin
      PGADMIN_CONFIG_SERVER_MODE: "False"
      PGADMIN_CONFIG_MASTER_PASSWORD_REQUIRED: "False"
    ports:
      - "5050:80"
    volumes:
      - ./pgadmin/servers.json:/pgadmin4/servers.json:ro
    depends_on:
      db:
        condition: service_healthy

volumes:
  pgdata:
\end{lstlisting}

\begin{itemize}
  \item \textbf{db} -- serverul PostgreSQL 17. În container rulează pe portul 5432, dar pe
        calculatorul gazdă este expus pe portul \textbf{5433} (configurabil prin \texttt{POSTGRES\_PORT}),
        deoarece portul 5432 era deja ocupat de o instalare locală PostgreSQL 18;
        volumul \texttt{pgdata} păstrează datele între reporniri; \emph{healthcheck}-ul
        verifică periodic, cu \texttt{pg\_isready}, dacă baza de date acceptă conexiuni.
  \item \textbf{pgadmin} -- interfață web pentru administrarea bazei de date, disponibilă la
        \url{http://localhost:5050}. Pornește doar după ce serviciul \texttt{db} este funcțional.
        Serverul bazei de date este înregistrat automat din fișierul \texttt{pgadmin/servers.json},
        iar cererea parolei principale (\emph{master password}) este dezactivată, fiind un mediu local.
\end{itemize}

\paragraph{Fișierul \texttt{.env}} conține configurările bazei de date, citite automat de
Docker Compose. Fișierul nu este urcat în Git (este listat în \texttt{.gitignore});
valorile sunt doar pentru mediul local de dezvoltare.
\begin{lstlisting}
POSTGRES_DB=second_chance_paws
POSTGRES_USER=scp
POSTGRES_PASSWORD=scp_dev_password
POSTGRES_PORT=5433
\end{lstlisting}

În repository este urcat fișierul \texttt{.env.example}, un model după care fiecare
dezvoltator își creează propriul \texttt{.env}.

\paragraph{Probleme întâlnite.}
\begin{itemize}
  \item \textbf{Port ocupat:} la prima pornire Docker a raportat
        \texttt{bind: address already in use} pentru portul 5432. Cauza: o instalare locală
        PostgreSQL 18. Soluția a fost maparea portului containerului pe 5433, fără a modifica
        instalarea existentă.
  \item \textbf{Serverul nu apărea în pgAdmin:} fișierul \texttt{servers.json} este importat doar
        la prima inițializare. Containerul a fost recreat împreună cu volumul său
        (\texttt{docker compose rm -sfv pgadmin}).
\end{itemize}

\paragraph{Verificare:}
\begin{lstlisting}
docker exec scp-db psql -U scp -d second_chance_paws -c "select version();"
\end{lstlisting}
Rezultat: \texttt{PostgreSQL 17.12 on aarch64-unknown-linux-musl}.

\paragraph{Pornire și oprire:}
\begin{lstlisting}
docker compose up -d     # porneste containerele in fundal
docker compose ps        # afiseaza starea containerelor
docker compose down      # opreste containerele (datele raman)
docker compose down -v   # opreste si sterge si datele
\end{lstlisting}

\subsection{Versionare cu Git și GitHub (\done)}
Codul este versionat cu Git și publicat într-un repository public:
\url{https://github.com/bonti23/second-chance-paws}.
Fișierul \texttt{.env} este exclus prin \texttt{.gitignore}. După fiecare pas finalizat
se face un \emph{commit} și un \emph{push}.

\section{Plan de dezvoltare}

\begin{center}
\begin{tabularx}{\textwidth}{@{}clX@{}}
\toprule
\textbf{Pas} & \textbf{Status} & \textbf{Conținut} \\
\midrule
1 & \done & PostgreSQL și pgAdmin în Docker \\
2 & \todo & Proiect Spring Boot conectat la baza de date \\
3 & \todo & Entități, migrări Flyway, date demonstrative \\
4 & \todo & Securitate (JWT, roluri), API cu paginare și filtre \\
5 & \todo & Frontend React și harta centrelor \\
6 & \todo & Docker Compose complet și integrarea AI \\
\bottomrule
\end{tabularx}
\end{center}

\end{document}
